\subsection{正规算子}

定义 4. --- 设 E 为希尔伯特空间，且 \(u \in \mathcal{L}(\mathrm{E})\)。若 \(u\) 与其伴随算子 \(u^{*}\) 交换，则称其为正规算子。

例如，设 \(u\) 是希尔伯特空间 \(E\) 的一个自同构，则它是正规的，因为我们有 \(uu^{*} = u^{*}u = 1_{E}\)。

命题 7. --- 为使 \(u \in \mathcal{L}(\mathrm{E})\) 是正规的，必要且充分的是：\(\| u(x) \| = \| u^{*}(x) \|\) 对所有 \(x \in \mathrm{E}\) 成立。

我们定义埃尔米特型 \(\Phi\)（在 \(E\) 上）如下

\[
\Phi (x, y) = \langle u u ^ {*} (x) | y \rangle - \langle u ^ {*} u (x) | y \rangle .
\]

为使 u 是正规的，必要且充分的是 \(\Phi = 0\)。由极化公式（V, p.~2），这等价于 \(\Phi(x, x) = 0\) 对所有 \(x \in E\) 成立。由于

\[
\Phi (x, x) = \left\| u ^ {*} (x) \right\| ^ {2} - \left\| u (x) \right\| ^ {2}.
\]

命题 8. --- 假设 \(u \in \mathcal{L}(\mathrm{E})\) 是正规的。设 \(\mathbf{N}\) 为 \(u\) 的核，\(\mathbf{M}\) 为 \(\mathbf{N}\) 在 \(\mathrm{E}\) 中的正交补；设 \(m\)、\(n\) 为满足 \(m + n \geqslant 1\) 的两个正整数。则 \(\mathbf{N}\) 是 \(u^{m}(u^{*})^{n}\) 的核，且 \(\mathbf{M}\) 既是 \(u^{m}(u^{*})^{n}\) 的初始子空间又是它的终子空间。特别地，\(\mathbf{M}\) 既是 \(u\) 与 \(u^{*}\) 的初始子空间，又是它们的终子空间，并且在 \(u\) 与 \(u^{*}\) 下稳定。

命题 7 表明 u 与 \(u^{*}\) 有相同的核 N。由 V, p.~41 的命题 4, (ii)，E 的子空间 M 在 u 与 \(u^{*}\) 下稳定，因为这对 \(N = M^{\circ}\) 成立；又因 \(M \cap N = \{0\}\)，u 与 \(u^{*}\) 在 M 上诱导的自同态都是单射。设 \(v = u^{m}(u^{*})^{n}\)；前面的论证表明 v 在 M（相应地 N）上的限制是单射（相应地为零），因而 N 是 v 的核。于是 \(M = N^{\circ}\) 是 v 的初始子空间。由 V, p.~41 的命题 4, (i)，v 的终子空间等于 \(v^{*}\) 的初始子空间。但 \(v^{*} = u^{n}(u^{*})^{m}\)，故由前所述，\(v^{*}\) 的初始子空间等于 M。

推论. --- 设 \(\lambda \in \mathbf{K}\)。E 的下列子空间相等：

\begin{enumerate}
\def\labelenumi{\alph{enumi})}
\item
  \(u\) 的关于 \(\lambda\) 的特征子空间；
\item
  \(u^{*}\) 的关于 \(\overline{\lambda}\) 的特征子空间；
\item
  \(u\) 的关于 \(\lambda\) 的准素子空间（换言之，按 LIE, VII, § 1, No.~1，即所有这样的向量 \(x\)（属于 \(E\)）组成的集：存在整数 \(n \geqslant 0\) 使 \((u - \lambda .1_{E})^{n}(x) = 0\)）；
\item
  \(u^{*}\) 的关于 \(\overline{\lambda}\) 的准素子空间。
\end{enumerate}

显然 \(w = u - \lambda \cdot 1_{E}\) 是 E 的正规自同态，因而由命题 8，E 的自同态 w、\(w^{*} = u^{*} - \overline{\lambda} \cdot 1_{E}\)、\(w^{n}\) 与 \((w^{*})^{n}\) 有相同的核。

